\begin{table}[ht!]
\centering
\caption{The decline read off the quarterly path, on contrasts that do not depend on where the adoption window is judged to begin.}
\label{tab:fixed_contrasts}
\footnotesize
\begin{tabular}{lcc}
\toprule
Contrast & Ages 22--25 & Ages 26--30 \\
\midrule
Calendar 2024 against calendar 2023 & $-0.0342^{*}$ (0.0084) & $-0.0235^{*}$ (0.0056) \\
Second half of 2024 against second half of 2023 & $-0.0449^{*}$ (0.0093) & $-0.0278^{*}$ (0.0058) \\
2024Q1 against 2023Q4, the quarters either side of the boundary & $+0.0041$ (0.0091) & $-0.0046$ (0.0045) \\
\bottomrule
\end{tabular}
\begin{minipage}{0.94\textwidth}\footnotesize\vspace{4pt}Linear combinations of the quarterly coefficients of Figure~3 of the paper, with standard errors $\sqrt{w'Vw}$ from that fit's own clustered covariance. Each row compares periods of equal length and identical calendar-quarter composition, so none of them depends on the January 2024 boundary, on the unequal lengths of the periods Equation~(2) partitions, or on the launch quarter. The first row is the boundary-free counterpart of the headline; it is not the headline re-estimated, since the pooled Poisson coefficient weights employer-months where these weight quarters equally. The last row is the pair of quarters either side of the boundary. $^{*}$ $p<0.05$. Source: script 84.
\end{minipage}
\end{table}
